\documentclass[10pt]{article}
\usepackage[a4paper,margin=0.65in]{geometry}
\usepackage{amsmath,amssymb,mathtools}
\usepackage[dvipsnames]{xcolor}
\usepackage{booktabs,array}
\usepackage[colorlinks=true,linkcolor=blue]{hyperref}
\numberwithin{equation}{section}
\setlength{\emergencystretch}{3em}\sloppy
\allowdisplaybreaks
\newcommand{\mf}[1]{\mathfrak{#1}}
\newcommand{\NLO}[1]{\frac{d^3N_{#1}}{d^3k}}
\title{Group II, Row IV from the amplitude: the BFKL term}
\date{}
\begin{document}
\maketitle
\noindent We start from your amplitude-level expression (with $\mf C$ and $\mf A^{(1)}$), do every step, and extract the $\log\frac{\vee}{k^+}$ term. Only BFKL is discussed. The result is the same as before:
\begin{equation*}
\boxed{\text{BFKL}\big|_{\rm II,IV}=\log\frac{\vee}{k^+}\,H_4}
\end{equation*}
This time it is obtained directly, without Fourier asymptotics. It also shows why $y=y'=z$.

\section{The expression and the conventions used}
Your expression:
\begin{multline}
\NLO{\rm NLO}=\frac{4}{(2\pi)^{3}}\int_{\Lambda}^{\vee-k^{+}}\frac{dp^{+}}{2\pi}\int_{z,w,w'}e^{-ik(w'-w)}\int_{y,y'}\\
\times\Big[\mf A^{\dagger(1)e'}_{k'}(P^+,-y')U^{\dagger e'c'}(y')U^{d'b}(z)\bar{\mf C}^{\dagger d'ac'}_{1jik'}-\bar{\mf A}^{\dagger(1)c'}_{k'}(P^+,-y')\bar{\mf C}^{\dagger d'ac'}_{1jik'}U^{d'b}(z)\Big]\\
\times\Big[\bar{\mf C}^{dac}_{1jik}U^{bd}(z)U^{ce}(y)\mf A^{(1)e}_{k}(P^+,-y)-U^{bd}(z)\bar{\mf C}^{dac}_{1jik}\bar{\mf A}^{(1)c}_{k}(P^+,-y)\Big],
\end{multline}
with $P^+\equiv p^++k^+$, $\bar{\mf C}_{1jik}\equiv\bar{\mf C}_{1jik}(P^+,y-w;p^+,w-z)$, and the conjugate one at $(y'-w';w'-z)$.

\paragraph{What each index is.}
\begin{itemize}
\item $\mf C^{bca}_{1jki}(k^+,z;p^+,x)$ is the splitting of a parent (momentum $k^+$, polarization $i$) into a gluon of momentum $p^+$ (polarization $j$) and a gluon of momentum $k^+-p^+$ (polarization $k$). The $\frac{k^+}{p^+}$ term has $j$ along $x$, so $j$ is the $p^+$ gluon.
\item In $\bar{\mf C}_{1jik}(P^+,y-w;p^+,w-z)$ the slots are relabelled:
\begin{itemize}
\item $k$ is the parent: momentum $P^+$, position $y$, colour $c$, contracted with $\mf A_k$;
\item $j$ is the unmeasured gluon: momentum $p^+$, position $z$, colour $d$;
\item $i$ is the measured gluon: momentum $k^+$, position $w$, colour $a$.
\end{itemize}
\end{itemize}
So, with $x\equiv w-z$, the tensor of $\bar{\mf C}_{1jik}$ is
\begin{equation}
T^{m}_{ijk}=\delta_{km}\delta_{ij}-\alpha\,\delta_{jm}\delta_{ik}-\beta\,\delta_{im}\delta_{jk},\qquad \alpha=\frac{P^+}{p^+},\quad\beta=\frac{P^+}{k^+}.
\end{equation}

\paragraph{Two things to check in the expression.}
\begin{enumerate}
\item In the conjugate bracket $U^{d'b}(z)$ should be $U^{\dagger d'b}(z)=U^{bd'}(z)$. Then $\sum_bU^{bd'}(z)U^{bd}(z)=\delta^{dd'}$: the Wilson lines of the unmeasured gluon cancel, as they must for an integrated-out gluon. I use this.
\item $\bar{\mf A}^{c}_k(P^+,-y)$ is $\mf A$ with $U(x)\rho(x)$: $\bar{\mf A}^{c}_k\propto\int_x\mathcal K^k(x-y)U^{ce}(x)\rho^e(x)$. Then the amplitude bracket is
\begin{equation}
\bar{\mf C}^{dac}_{1jik}U^{bd}(z)\,\frac{ig}{\sqrt2\pi\sqrt{P^+}}\,L^{c}_k(y),\qquad
L^{c}_k(y)\equiv\int_x\mathcal K^k(x-y)\big[U^{ce}(y)-U^{ce}(x)\big]\rho^e(x),
\end{equation}
which is the LO amplitude of the parent gluon at $y$. The overall sign convention of $\mf A$ cancels between amplitude and conjugate.
\end{enumerate}

\section{Step 1: the $\delta$-function puts the parent at $y=\xi w+(1-\xi)z$}
$\mf C(k^+,z;p^+,x)\propto\delta^{(2)}\big(z+\frac{p^+}{k^+}x\big)$. In $\bar{\mf C}(P^+,y-w;p^+,w-z)$ this is
\begin{equation}
\delta^{(2)}\Big(y-w+\frac{p^+}{P^+}(w-z)\Big)\quad\Rightarrow\quad y=\frac{k^+}{P^+}\,w+\frac{p^+}{P^+}\,z=\xi\,w+(1-\xi)\,z,\qquad \xi\equiv\frac{k^+}{p^++k^+}.
\end{equation}
The parent sits at the $+$-momentum-weighted centre of its two daughters. Likewise $y'=\xi w'+(1-\xi)z$. The $y,y'$ integrals are done by these $\delta$'s. Hence
\begin{equation}
y=z+\xi(w-z),\qquad y'=z+\xi(w'-z),\qquad y'-y=\xi(w'-w).\label{eq:y}
\end{equation}

\section{Step 2: colour, normalization and the $p^+$ integral}
\paragraph{Colour.} With $\delta^{dd'}$ from Section 1, $\sum_{a,d}f^{c'da}f^{cda}=N_c\,\delta^{cc'}$. So amplitude and conjugate have the same parent colour $c$, contracted between $L^c_k(y)$ and $L^{c}_{k'}(y')^\dagger$.

\paragraph{Factors.}
\begin{itemize}
\item $|\mf C|^2\to\frac{g^2N_c}{4\pi^2\,2P^+}\frac{p^+k^+}{P^{+2}}\,\frac{x^m x'^{m'}}{x^2x'^2}\,T^m_{ijk}T^{m'}_{ijk'}$, with $x=w-z$, $x'=w'-z$.
\item $|\mf A|^2\to\frac{g^2}{2\pi^2P^+}L_kL'_{k'}$.
\item $\frac{p^+k^+}{P^{+2}}=\xi(1-\xi)$.
\end{itemize}
\paragraph{Change of variable.} $p^+=k^+\frac{1-\xi}{\xi}$, $dp^+=-\frac{k^+}{\xi^2}d\xi$, so $\int_\Lambda^{\vee-k^+}\frac{dp^+}{P^{+2}}=\frac{1}{k^+}\int_a^{\xi_{\max}}d\xi$, with $a=\frac{k^+}{\vee}$ and $\xi_{\max}=\frac{k^+}{k^++\Lambda}$.
\paragraph{Normalization.} $\frac{4}{(2\pi)^3}\cdot\frac{1}{2\pi k^+}\cdot\frac{g^4N_c}{16\pi^4}=\frac{1}{(2\pi)^3}\frac{g^4N_c}{8\pi^5k^+}\equiv\mathcal N$.

\section{Step 3: the polarization sum (any $d_\perp$)}
\begin{equation}
\sum_{i,j}T^{m}_{ijk}T^{m'}_{ijk'}=(\alpha^2+\beta^2)\,\delta_{mm'}\delta_{kk'}+2\alpha\beta\,\delta_{mk'}\delta_{km'}+(d_\perp-2\alpha-2\beta)\,\delta_{km}\delta_{k'm'} .
\end{equation}
Multiply by $\xi(1-\xi)$, using $\alpha=\frac{1}{1-\xi}$ and $\beta=\frac1\xi$:
\begin{equation}
n_1=\underbrace{\frac{1-\xi}{\xi}}_{\beta^2\xi(1-\xi)}+\frac{\xi}{1-\xi},\qquad n_2=2,\qquad n_3=d_\perp\,\xi(1-\xi)-2 .
\end{equation}
(Checked for $d=2,3,4$ and symbolically: \texttt{from\_amplitude/checks.py}.)

\section{Step 4: Row IV, exact}
\begin{multline}
\text{Row IV}=\mathcal N\int_{a}^{\xi_{\max}}d\xi\int_{z,w,w'}e^{-ik(w'-w)}\Big[n_1\,\mathcal K(w-z)\cdot\mathcal K(w'-z)\;L(y)\cdot L^\dagger(y')\\
+n_2\,\big(\mathcal K(w-z)\cdot L^\dagger(y')\big)\big(\mathcal K(w'-z)\cdot L(y)\big)+n_3\,\big(\mathcal K(w-z)\cdot L(y)\big)\big(\mathcal K(w'-z)\cdot L^\dagger(y')\big)\Big],\label{eq:exact}
\end{multline}
with $y,y'$ from Eq.~\eqref{eq:y} and the colour $c$ summed in each product of $L$'s.
\begin{itemize}
\item $\vee$ appears only in the lower limit $a=\frac{k^+}{\vee}$.
\item There is no explicit $\log(\frac{\vee}{k^+}-1)$ here. The one in line 2 of the three-line form is produced by the rewriting into a plus prescription and cancels again (see \texttt{RowIV\_BFKL\_only.pdf}).
\end{itemize}

\section{Step 5: the $\log\frac{\vee}{k^+}$}
In Eq.~\eqref{eq:exact} the only $\xi$-dependence, apart from $n_{1,2,3}$, is through $y=z+\xi(w-z)$ and $y'=z+\xi(w'-z)$. As $\xi\to0$ these go smoothly to $y,y'\to z$.
\begin{itemize}
\item $n_2=2$ and $n_3=d_\perp\xi(1-\xi)-2$ are finite at $\xi=0$. So their integrals have no $\log\frac1a$.
\item $\frac{\xi}{1-\xi}$ in $n_1$ is finite at $\xi=0$ (it is the $\xi\to1$, i.e.\ JIMWLK, end).
\item \textbf{Only $\beta^2\xi(1-\xi)=\frac{1-\xi}{\xi}$ is singular}, and $\int_a\frac{d\xi}{\xi}=\log\frac{\vee}{k^+}+\dots$.
\end{itemize}
So the coefficient of $\log\frac{\vee}{k^+}$ is the $n_1$ integrand at $\xi=0$, i.e.\ at $y=y'=z$:
\begin{equation}
\text{Row IV}=\log\frac{\vee}{k^+}\;\mathcal N\int_{z,w,w'}e^{-ik(w'-w)}\,\mathcal K(w-z)\cdot\mathcal K(w'-z)\;L(z)\cdot L^\dagger(z)\;+\;\big(\text{no }\log\vee\big).
\end{equation}
\begin{itemize}
\item At $\xi=0$ the $w,w'$ integrals are finite: $\int_{w,w'}e^{-ik(w'-w)}\mathcal K(w-z)\cdot\mathcal K(w'-z)=\frac{4\pi^2}{\mathbf k^2}$, with no $\log\mu$ and no pole.
\item At finite $\xi$ Eq.~\eqref{eq:exact} has the collinear pole. Its coefficient, however, is the Fourier transform of a smooth function at momentum $\mathbf k/\xi$. It vanishes faster than any power as $\xi\to0$, so it does not touch the $\log\vee$.
\end{itemize}

\paragraph{Why $y=y'=z$.} This answers your question. $y$ and $y'$ are the parent's positions in the amplitude and in the conjugate. By Eq.~\eqref{eq:y}, $y=\xi w+(1-\xi)z$: the parent is at the $+$-momentum-weighted centre of its daughters. In the BFKL region $\xi\to0$ ($p^+\gg k^+$), the unmeasured gluon at $z$ carries all the $+$ momentum. So the parent sits at $z$, in the amplitude and in the conjugate alike. The same point $z$ was found earlier from the Fourier transform in the $(y,y')$ variables; here it is just kinematics.

\section{Result}
Write out $L(z)\cdot L^\dagger(z)$:
\begin{equation}
\boxed{\begin{aligned}
\text{BFKL}\big|_{\rm II,IV}={}&\log\frac{\vee}{k^+}\,H_4,\\
H_4={}&\frac{1}{(2\pi)^{3}}\frac{g^{4}N_c}{8\pi^{5}}\frac{1}{k^{+}}\int_{x,x',z,w,w'}e^{-ik(w'-w)}\,\mathcal K(w-z)\cdot\mathcal K(w'-z)\;\mathcal K(x-z)\cdot\mathcal K(x'-z)\\
&\times\big[U(z)-U(x')\big]^{e'c}\big[U(z)-U(x)\big]^{ce}\,\rho^{e'}(x')\rho^{e}(x).
\end{aligned}}
\end{equation}
\begin{itemize}
\item This is exactly the $H_4$ obtained before from the three-line form, with the same prefactor $\frac{g^4N_c}{8\pi^5}$, the same kernels and the same Wilson lines.
\item Expanded, the bracket gives your four lines: $U(x')U(x)\rho\rho$, $\rho\rho$ (from $U(z)U^\dagger(z)=1$), $-U(z)U(x)\rho\rho$, $-U(x')U(z)\rho\rho$.
\item The colour orientation follows your $\mf A$, $\bar{\mf A}$; in your LO notation the first factor reads $U^{e'c}(y')$ and the second $U^{ec}(y)$.
\end{itemize}
\paragraph{Physical picture.} The term comes from $\beta^2=(P^+/k^+)^2$ in $|\mf C|^2$. That is the piece where the \emph{measured} gluon's polarization is along $x=w-z$ ($\delta_{im}$) and the parent's polarization passes to the fast gluon ($\delta_{jk}$). It is the soft emission of the measured gluon, with kernel $\mathcal K(w-z)$, from a fast gluon produced at $z$ with the LO amplitude $L(z)$. Its weight $\frac{1-\xi}{\xi}$ is the $\frac1\xi$ pole of $P_{gg}$.

\section{Comparison with your three-line form}
To compare, go to the variables $(y,y')$: by Eq.~\eqref{eq:y}, $d^2w\,d^2w'=\xi^{-4}d^2y\,d^2y'$ and $\frac{x^m}{x^2}=\xi\frac{(y-z)^m}{(y-z)^2}$. This gives the $\frac{d\xi}{\xi^2}$ and the phase $e^{-ik(y'-y)/\xi}$. The $z$ integral is then
\begin{equation}
\int_z\frac{(y-z)^m(y'-z)^{m'}}{(y-z)^2(y'-z)^2}=A\,\delta_{mm'}+B'\,\frac{X^mX^{m'}}{X^2},\qquad B'=(d_\perp-2)A\to-\pi .
\end{equation}
$A$ carries the pole and the $\ln X^2$; the finite $B'=-\pi$ is checked numerically. Then:
\begin{itemize}
\item The $L\cdot L'$ coefficient is $n_1(2d_\perp-2)A+d_\perp\xi(1-\xi)A=d_\perp A\big[\xi(1-\xi)+\frac{1-\xi}{\xi}+\frac{\xi}{1-\xi}\big]+(d_\perp-2)A\,n_1$. This is your bracket in line 1.
\item The $\frac{(X\cdot L)(X\cdot L')}{X^2}$ coefficient is $d_\perp\xi(1-\xi)B'/\xi^2=-\pi\,d_\perp\frac{1-\xi}{\xi}$. This is your $-d_\perp\frac{1-\xi}{\xi}$ term in line 3.
\item The $\frac{\xi}{1-\xi}$ term with the upper limit $\xi_{\max}$ gives your $\log\frac{k^+}{\Lambda}$ line.
\end{itemize}
So the $\xi$-structure of your three lines is reproduced. I did not recheck the $\epsilon$-constants of $A$ here; they do not enter the BFKL term. Both routes give the same $\log\frac{\vee}{k^+}H_4$.

\paragraph{Checks} (\texttt{from\_amplitude/checks.py}):
\begin{itemize}
\item polarization sum, exact for $d=2,3,4$;
\item $n_{1,2,3}$, symbolic;
\item $f^{c'da}f^{cda}=N_c\delta^{cc'}$ for SU(3), to $10^{-15}$;
\item normalization $\mathcal N$, symbolic;
\item $B'=-3.14159265359$.
\end{itemize}
\end{document}
